%---------------------------------------------
\section{Experiments}
\label{experisment}
%---------------------------------------------

\noindent\textbf{Setting up.} We conduct experiments based on an open-source T2V diffusion model VideoCrafter~\citep{chen2023videocrafter1} for both singe-prompt and multi-prompt longer video generations. The video diffusion model is trained on $16$ frames and is required to sample $64$ frames in the inference stage. The window and stride size are set to $U=16$, $S=4$ as default.

\noindent\textbf{Evaluation Metrics}. To evaluate our paradigm, we report Frechet Video Distance (FVD) \citep{unterthiner2018towards}, Kernel Video Distance (KVD) \citep{kvd} and Clip Similarity (CLIP-SIM) \citep{radford2021learning}. Since the longer inference methods are supposed to keep the quality of the original fixed-length inference, we calculate the FVD between original generated short videos and subset generated longer videos with corresponding lengths. CLIP-SIM is used to measure the content consistency of generated videos by calculating the semantic similarity among adjacent frames of generated videos.
%


\begin{table}[t]
\caption{Quantitative comparison on longer video generation.}\vspace{-0.5em}
\label{tbl:score}
\begin{center}
\scalebox{0.8}{\begin{tabular}{lcccc}
\hline
\noalign{\smallskip}
\multicolumn{1}{c}{\bf Method}  &\multicolumn{1}{c}{\bf FVD ()} &\multicolumn{1}{c}{\bf KVD ()} &\multicolumn{1}{c}{\bf CLIP-SIM ($\uparrow$)} &\multicolumn{1}{c}{\bf Inference Time ()} \\
\noalign{\smallskip}
\hline
\noalign{\smallskip}
Direct       &737.61  &359.11 & 0.9104 & \textbf{21.97s} \\
Sliding      &224.55  &44.09 & 0.9438 & 36.76s \\
GenL         &177.63  &21.06 & 0.9370 & 77.89s \\
Ours         &\textbf{85.83}  &\textbf{7.06} & \textbf{0.9732} & 25.75s \\
\noalign{\smallskip}
\hline
\end{tabular}}\vspace{-1.0em}
\end{center}
\end{table}

\begin{figure*} %[t]
%
\centering
\includegraphics[width=0.98\linewidth]{Figures/fig_comp_long.pdf}\vspace{-1.3em}
\caption{Qualitative comparisons of longer video generation. Left prompt: ``\textit{A chihuahua in astronaut suit floating in space, cinematic lighting, glow effect}''. Right prompt: ``\textit{A very happy fuzzy panda dressed as a chef eating pizza in the New York street food truck}''.
}
\vspace{-1.0em}
\label{fig:comp_long}
\end{figure*}

\subsection{Longer Video Generation}

We mainly compare our proposed FreeNoise to other tuning-free longer video generation methods with diffusion models. We first directly sample $64$ frames (\textbf{Direct}). Then we adopt temporal sliding windows so that the temporal attention module can always process a fixed number of frames (\textbf{Sliding}). The closest work to our paradigm is Gen-L-Video (\textbf{GenL}), which extends the video smoothly by merging some overlapping sub-segments during the denoising process.

The synthesis results are shown in Figure~\ref{fig:comp_long}. In the first line, the dog has severe artifacts and the background of space is not clear. Obviously, directly sampling $64$ frames through a model trained on $16$ frames will bring poor quality results due to the training-inference gap. When we use temporal sliding windows, the training-inference gap is eliminated thus more vivid videos are generated. However, this operation ignores the long-range visual consistency thus the resulting subject and background both look significantly different among different frames. 
%
Gen-L-Video promotes the integration of frames by averaging the overlapping sub-segments and performs better in some cases. However, it fails to maintain long-range visual consistency and suffers from content mutation. Benefiting from noise rescheduling, all sub-segments in our paradigm share similar main subjects and scenarios while still containing considerable content variation, keeping our main content even when the generated video becomes longer.
%
Results shown in Figure~\ref{fig:comp_long} exhibit that our FreeNoise successfully renders high-fidelity longer videos, outperforming all other methods.

In addition, we also compare the operation time of those methods on NVIDIA A100. As presented in Table~\ref{tbl:score}, it is observed that Gen-L-Video exhibits the longest inference time, nearly four times longer than direct inference. This is primarily attributed to its default setting, which involves the nearly global sampling of the entire set of latents four times. However, our paradigm only brings less than $20\%$ extra inference time by limiting most additional calculations within the temporal attention layers.

Table ~\ref{tbl:score} shows quantitative results. The quality of videos generated by direct inference is extremely damaged by the training-inference gap, obtaining the worst FVD and KVD. The video quality from the sliding method and Gen-L-Video is obviously improved but still worse than the results generated by FreeNoise. Our FreeNoise also gains the best CLIP-SIM, indicating the superiority of our method in content consistency.

In addition, we conducted a user study to evaluate our results by human subjective perception. Users are asked to watch the generated videos of all the methods, where each example is displayed in a random order to avoid bias, and then pick up the best one in three evaluation aspects. As shown in Table ~\ref{tbl:user}, our approach achieves the highest scores for all aspects: content consistency, video quality, and video-text alignment, outperforming baseline methods by a large margin. Especially for content consistency, our method has received almost twice as many votes as the second place.

\begin{table}[t]
\caption{User study. Users are required to pick the best one among our proposed FreeNoise with the other baseline methods in terms of content consistency, video quality, and video-text alignment.}\vspace{-0.5em}
% 
\label{tbl:user}
\begin{center}
\scalebox{0.8}{\begin{tabular}{lccc}
\hline
\noalign{\smallskip}
\multicolumn{1}{c}{\bf Method}  &\multicolumn{1}{c}{\bf Content Consistency} &\multicolumn{1}{c}{\bf Video Quality}
&\multicolumn{1}{c}{\bf Video-Text Alignment} \\
\noalign{\smallskip}
\hline
\noalign{\smallskip}
Direct       &11.73\%  &10.80\%  &11.11\% \\
Sliding   &6.17\%  &6.79\%  &8.02\% \\
GenL            &24.38\%  &26.85\%  &29.63\% \\
Ours                   &\textbf{57.72\%}  &\textbf{55.56\%}  &\textbf{51.23\%} \\
\noalign{\smallskip}
\hline
\end{tabular}}\vspace{-1.0em}
\end{center}
\end{table}

%
\begin{figure*} %[t]
%
\centering
\includegraphics[width=0.98\linewidth]{Figures/fig_comp_multi.pdf}\vspace{-1.3em}
\caption{Qualitative comparisons of multi-prompt video generation. Left multi-prompt: ``\textit{A camel running on the snow field}'' $\to$ ``\textit{A camel standing on the snow field}''. Right multi-prompt: ``\textit{An astronaut resting on a horse}'' $\to$ ``\textit{An astronaut riding a horse}''.
}
\vspace{-1.0em}
\label{fig:comp_multi}
\end{figure*}


\paragraph{Multi-prompt Video Generation.}
We extend our paradigm for multi-prompt video generation by introducing the Motion Injection method. As shown in Figure~\ref{fig:comp_multi}, our method achieves coherent visual coherence and motion continuity: The camel gradually changes from running to standing while the distant mountains remain consistent appearances; The astronaut changes from resting on a horse to riding a horse naturally. However, when we purely use the strategy of noise rescheduling without motion injection, the scene will undergo unexpected changes because a new prompt often introduces unexpected new contents other than the text description due to the inherent properties of the Stable Diffusion model. But it can still work in some cases when the main objects and scenarios are not obviously changed by the new prompt (like the bigfoot case in Figure~\ref{fig:comp_multi}).
%
We also compare with the existing tuning-free state-of-the-art method Gen-L-Video. Figure~\ref{fig:comp_multi} shows that Gen-L-Video also achieves the conversion of two actions. However, due to the drawbacks of content mutation, its generated objects and scenarios are meaninglessly changed over time.

\subsection{Ablation Study}

\textbf{Ablation for Noise Rescheduling.} As noise rescheduling plays an essential role in our method, we typically conduct an ablation to validate its importance, namely removing it from our proposed inference paradigm. In addition, we also implement another variant of our method with the local noise shuffle unit size as $S=8$ (the sliding window stride is also changed to $8$ accordingly).

As shown in Figure~\ref{fig:abl_long}, without noise rescheduling, our method fails to keep content consistent. Although each frame still matches the text description, they are not semantically connected.
%
And when the sliding window stride is $8$, the synthesized features across windows are interacted in a less tight manner. For example, the shape of the bowl is changed gradually in Figure~\ref{fig:abl_long}. 
%
Since the stride value of $4$ is able to achieve enough content consistency, we do not consider the smaller stride value of $2$, which will bring the double extra inference time.


\begin{figure*} %[t]
%
\centering
\includegraphics[width=0.98\linewidth]{Figures/fig_abl_long.pdf}\vspace{-1.3em}
\caption{Ablation study on noise rescheduling. Left prompt: ``\textit{A cat eating food out of a bowl}''. Right prompt: ``\textit{A video of milk pouring over strawberries, blueberries, and blackberries}''.
}\vspace{-1.0em}
\label{fig:abl_long}
\end{figure*}

\begin{wrapfigure}{r}{0.45\textwidth}
    \centering
    \vspace{-0.5em}
    \includegraphics[width=0.45\textwidth]{Figures/fig_abl_multi.pdf}\vspace{-1.0em}
    \caption{Ablation study on motion injection. The multi-prompt is: ``\textit{An astronaut riding a horse}'' $\to$ ``\textit{An astronaut resting on a horse}''.
    }\vspace{-0.5em}
    \label{fig:abl_multi}
\end{wrapfigure}
%
\textbf{Ablation for Motion Injection.}
To show the effectiveness of our design choices in Motion Injection, we study on two main hyper-parameters---layer selection and timestep selection, as expressed in Equation~\ref{eq:mj_1}. In our design, the last $L$ cross-attention layers of the U-Net (e.g. the decoder part) will always be provided with the target prompt $\widetilde{P}$ across all the denoising steps and $P_1$ only maintains the layout and visual details through the cross-attention layers before the $L$ in some denoising steps. For comparison, we construct another two variations that $P_1$ control the only decoder part and all layers, respectively. Besides, the third variation allows $P_1$ to control the only decoder part across all the denoising steps.
%
Figure~\ref{fig:abl_multi} shows the results of those variations. When $P_1$ is only allowed to control the decoder part, the running motion of the horse is suppressed. Meanwhile, the appearance of the horse is changed obviously (shown in Figure~\ref{fig:abl_multi}(a)). When $P_1$ is only allowed to control both the encoder and the decoder part, the appearance of the horse is kept but the running motion is still suppressed (shown in Figure~\ref{fig:abl_multi}(b)). It is because the decoder features are more tightly aligned with the semantic structures, as observed in MasaCtrl~\citep{cao2023masactrl}. 
%
Compared to layer factors, the choice of timesteps has relatively less influence. Even if we allow $P_1$ to control the only decoder part across all the denoising steps, the horse can still switch smoothly from running to standing. However, this behavior influences the generation of layout and visual details. As a result, a section of the horse's leg will suddenly disappear (shown in Figure~\ref{fig:abl_multi}(c)), while this missing leg appears bent and curled up in the same frame of our final result (shown in Figure~\ref{fig:abl_multi}(d)). Therefore, the selection of time steps can help to achieve more precise control over the content level.
%
